\documentclass[12pt,a4paper]{article}

\usepackage[margin=1in]{geometry}
\usepackage{graphicx}
\usepackage{xcolor}
\usepackage{listings}
\usepackage{hyperref}
\usepackage{titlesec}
\usepackage{fancyhdr}

% ── Page Style ──
\pagestyle{fancy}
\fancyhf{}
\rhead{CSE341 -- Internet Programming}
\lhead{Lab 3}
\cfoot{\thepage}

% ── Code Listing Style ──
\lstset{
  basicstyle=\ttfamily\small,
  backgroundcolor=\color{gray!10},
  frame=single,
  rulecolor=\color{gray!40},
  breaklines=true,
  numbers=left,
  numberstyle=\tiny\color{gray},
  tabsize=2,
  showstringspaces=false,
  keywordstyle=\color{blue!70!black}\bfseries,
  commentstyle=\color{green!50!black},
  stringstyle=\color{red!60!black}
}

\begin{document}

% ── Title Page ──
\begin{titlepage}
  \centering
  \vspace*{2cm}
  {\Huge\bfseries My Interactive Calculator Webpage\par}
  \vspace{1cm}
  {\Large CSE341 -- Internet Programming\par}
  {\Large Lab 3\par}
  \vspace{2cm}
  {\large
    \textbf{Name:} Seif Elden \\[6pt]
    \textbf{Date:} \today
  }
  \vfill
  {\large Ain Shams University \\ Faculty of Engineering}
\end{titlepage}

% ── Table of Contents ──
\tableofcontents
\newpage

% ════════════════════════════════════════════
\section{Objective}
Build an interactive calculator webpage using HTML, CSS, and JavaScript that performs basic mathematical operations (addition, subtraction, multiplication, division, and modulus).

% ════════════════════════════════════════════
\section{Project Structure}
The project consists of three files, each handling a separate concern:

\begin{itemize}
  \item \textbf{index.html} -- Document structure and calculator layout.
  \item \textbf{style.css} -- Visual styling with a glassmorphism design.
  \item \textbf{script.js} -- Calculator logic, DOM manipulation, and event handling.
\end{itemize}

% ════════════════════════════════════════════
\section{Implementation Details}

\subsection{HTML (index.html)}
The HTML file defines a semantic structure for the calculator:

\begin{itemize}
  \item A \texttt{.display} section with two elements: one for the running expression and one for the current result.
  \item A \texttt{.buttons} grid containing number buttons (0--9), operator buttons (+, $-$, $\times$, $\div$, \%), a decimal point, clear (AC), backspace ($\leftarrow$), and equals (=).
  \item Each button uses a \texttt{data-action} attribute to identify its function, keeping the HTML clean and the JS decoupled from the display text.
\end{itemize}

\subsection{CSS (style.css)}
Key creative design choices:

\begin{itemize}
  \item \textbf{Glassmorphism:} The calculator uses \texttt{backdrop-filter: blur()} and semi-transparent backgrounds to achieve a frosted-glass look.
  \item \textbf{Gradient Background:} A deep purple-to-dark gradient gives the page a modern aesthetic.
  \item \textbf{Grid Layout:} CSS Grid arranges the buttons in a clean 4-column layout.
  \item \textbf{Micro-interactions:} Buttons scale down on click and the result text pops on update, providing tactile feedback.
  \item \textbf{Responsive:} A media query adapts sizing for screens under 400px.
\end{itemize}

\subsection{JavaScript (script.js)}
The JS follows the DOM $\rightarrow$ Event $\rightarrow$ Logic pattern:

\begin{enumerate}
  \item \textbf{DOM Selection:} \texttt{getElementById} and \texttt{querySelectorAll} select the display and button nodes.
  \item \textbf{Event Listeners:} Click events on each button and a \texttt{keydown} listener for full keyboard support.
  \item \textbf{Logic Functions:}
    \begin{itemize}
      \item \texttt{appendNumber()} -- builds the current operand string.
      \item \texttt{appendOperator()} -- appends the operand and operator to the expression.
      \item \texttt{handleBackspace()} -- removes the last character or operator.
      \item \texttt{calculate()} -- evaluates the full expression safely and displays the result.
      \item \texttt{resetCalc()} -- clears all state.
      \item \texttt{render()} -- updates the DOM display elements.
    \end{itemize}
\end{enumerate}

% ════════════════════════════════════════════
\section{Screenshots}

\subsection{Calculator -- Initial State}
% ┌──────────────────────────────────────────┐
% │  Replace "screenshot1.png" with your     │
% │  actual screenshot file name.            │
% └──────────────────────────────────────────┘
\begin{figure}[h!]
  \centering
  \fbox{\parbox{0.7\textwidth}{\vspace{5cm}\centering\textit{Screenshot 1: Place your screenshot here\\(replace with \textbackslash includegraphics)}\vspace{5cm}}}
  % Uncomment the line below and comment the \fbox block above after adding your screenshot:
  % \includegraphics[width=0.7\textwidth]{screenshot1.png}
  \caption{Calculator in its initial state showing the default display.}
\end{figure}

\subsection{Calculator -- After Performing a Calculation}
% ┌──────────────────────────────────────────┐
% │  Replace "screenshot2.png" with your     │
% │  actual screenshot file name.            │
% └──────────────────────────────────────────┘
\begin{figure}[h!]
  \centering
  \fbox{\parbox{0.7\textwidth}{\vspace{5cm}\centering\textit{Screenshot 2: Place your screenshot here\\(replace with \textbackslash includegraphics)}\vspace{5cm}}}
  % Uncomment the line below and comment the \fbox block above after adding your screenshot:
  % \includegraphics[width=0.7\textwidth]{screenshot2.png}
  \caption{Calculator displaying the result of a mathematical operation.}
\end{figure}

% ════════════════════════════════════════════
\section{Conclusion}
This lab demonstrated building a fully functional calculator webpage using the separation-of-concerns principle: HTML for structure, CSS for presentation, and JavaScript for behavior. The JS code was organized around the DOM--Event--Logic pattern, making it maintainable and efficient. The glassmorphism CSS design gives the calculator a polished, modern look.

\end{document}
